\subsection*{Biological agents used}\label{sec:agents}
Organisms and viruses, both henceforth collectively referred to as ‘organisms,’ assigned to risk group (RG) 2 may cause infectious diseases to humans. According to the current scientific knowledge, the risk of working with RG 2 organisms low to medium. A risk assessment carried out by the Health \& Safety Office (HSO) becomes necessary when genetic engineering experiments involving RG 2 organisms as gene donor or recipient are planned. If the resulting genetically modified organisms (GMO) are assessed to belong to RG 2, handling them must be restricted to safety level 2 (S\,2) laboratories. 
\newline

In the “Structures” S 2 facility, work is carried out with both genetically modified and wildtype \emph{Mycoplasma pneumoniae} and genome libraries of \emph{Streptococcus salivarius}, \emph{Bacteroides uniformus}, and \\emph{Bacteroides fragilis}. \emph{Salmonella typhimurium} is modified to do the heterologous expression of different fluorescent proteins. Further, Mumps orthorubula and Sinbis virus are used in conjunction with genetically modified HeLa cells. The infected cells are grown and transfected with mammalian expression vectors to express target proteins fused to fluorescent proteins. The viruses can proliferate in the infected cells and have the potential to spread. Further, GMO variants of the Sindbis virus, expressing fluorescent proteins, are used for infection studies. Replication-deficient lentiviral particles are used to insert genes encoding for fluorescent proteins into the genome of mammalian cell lines.
\newline

In most cases, infections with \textbf{\emph{Mycoplasma pneumoniae}} effect the upper respiratory tract causing fever, worsening dry cough, malaise, and headache. Further symptoms are chilliness, scratchy sore throat, sore chest and tracheal tenderness as a result of the protracted cough and in rare cases pleuritic chest pain. Radiologically confirmed pneumonia develops in 5-10 \% of cases. Only rarely extrapulmonary syndromes, including cardiologic, neurologic, and dermatologic findings have been described. In almost all patients, the pneumonia resolves without any serious complications. Infections with \emph{M. pneumoniae} are transmitted mainly by droplet inhalation or direct contact with infected persons. The incubation period averages three weeks, in contrast to that of influenza or other viral pneumonias, which generally is a few days. 
\newline

\textbf{\emph{Mycobacterium smegmatis}}, a rapidly growing non-tuberculous mycobacterium, is widely used as a surrogate model for pathogenic species due to its genetic tractability and biosafety profile. Although generally regarded as nonpathogenic, it has been implicated in opportunistic infections, particularly in immunocompromised hosts or following trauma and surgical interventions. Documented cases include cutaneous and soft tissue infections, osteomyelitis, catheter-related bacteremia, and, less frequently, pulmonary involvement. Pathogenesis is linked to its intrinsic resistance to many antibiotics, robust biofilm formation, and survival within macrophages. Laboratory exposure risk remains low under biosafety level 2 precautions; however, breaches in aseptic technique or contact with open wounds can facilitate transmission. Treatment may require prolonged multidrug regimens, often including macrolides, aminoglycosides, or fluoroquinolones, guided by susceptibility testing. 
\newline

Mumps is the disease caused by the \textbf{mumps virus}. Initial symptoms are non-specific and include fever, headache, malaise, muscle pain, and loss of appetite. These symptoms are usually followed by painful swelling of the parotid glands, called parotitis, which is the most common symptom of infection. Symptoms typically occur 16 to 18 days after exposure to the virus and resolve within two weeks. About one third of infections are asymptomatic. Complications include deafness and a wide range of inflammatory conditions, of which inflammation of the testes, breasts, ovaries, pancreas, meninges, and brain are the most common. Testicular inflammation may result in reduced fertility and, rarely, sterility. 
\newline

\textbf{Sindbis virus} infections in humans typically cause a mild illness known as Sindbis fever. Common symptoms include fever, rash, and joint pain (arthralgia), particularly in the wrists, knees, and ankles. Some individuals may also experience muscle pain, fatigue, headache, and swollen lymph nodes. The rash usually appears on the trunk and limbs. Symptoms generally develop a few days after being bitten by an infected mosquito and can last for several weeks, especially joint pain. Most cases resolve without complications, but persistent arthritis may occur in some individuals. Severe cases are extremely rare. The infection is not contagious between people.
\newline

The genus \textbf{\emph{Bacteroides}} comprises a diverse group of Gram-negative bacteria that are commonly found in the human gut microbiome. \emph{Bacteroides} species are known to produce a range of virulence factors, including lipopolysaccharides, proteases, and capsular polysaccharides, which can damage host tissues and promote inflammation. In addition, \{emph{Bacteroides} species are often involved in polymicrobial infections, where they can interact with other microorganisms to exacerbate disease. B. fragilis, in particular, is a well-known pathogen that is responsible for a range of infections, including intra-abdominal infections, bacteremia, and sepsis. This bacterium produces several virulence factors, including a toxin called ’\emph{Bacteroides fragilis} toxin’, which can cause damage to the intestinal epithelium and promote inflammation. Furthermore, the emergence of antibiotic-resistant strains of \emph{Bacteroides} has become a growing concern. These resistant strains can be difficult to treat and can lead to severe and potentially life-threatening infections.
\newline

\textbf{\emph{Streptococcus salivarius}} is a Gram-positive bacterium that is commonly found in the human oral cavity. While this bacterium is generally considered to be non-pathogenic, there have been instances where it has been implicated in infectious diseases. Studies have shown that certain strains of \emph{Streptococcus salivarius} can produce virulence factors, such as bacteriocins and adhesins, that can damage host tissues and promote inflammation. In addition, \emph{Streptococcus salivarius} has been associated with several infectious diseases, including bacterial endocarditis, meningitis, and sepsis. As with \emph{Bacteroides} sp., the emergence of antibiotic-resistant strains is a reason for particular concern.
\newline

\textbf{\emph{Salmonella typhimurium}} may cause salmonellosis, a diarrhoeal disease. As it is associated with animals (rather than through contact with infected people), it is usually a foodborne disease, spread through consumption of contaminated foods, especially eggs and poultry. A fair number of cases of transmission have also been attributed to exposure to pets infected with the bacterium. Symptoms appear within approximately 8 to 72 hours after ingesting the contaminated food and may last 2 to 7 days. They include diarrhea (usually non-bloody), nausea, vomiting, abdominal pain, headache, and fever.
\newline

\textbf{Lentivirus} (lenti-, Latin for \emph{slow}) is a genus of slow viruses of the Retroviridae family, which is characterised by a long incubation period. Lentiviruses can deliver a significant amount of genetic information into the DNA of the host cell and hence constitute one of the most efficient gene delivery vectors. Examples of lentiviruses are HIV$^a$, SIV$^a$, and FIV. The lab strains of lentiviruses used for genetic manipulation lack the capability for replication and do not cause infectious diseases. However, they pose a weak risk of integrating into the lab workers  DNA and causing cancer if droplets containing the lentiviruses are inhaled.\\
\footnotesize a) the wild type viruses are assigned to risk group 3$^{**}$)\normalsize

